% CHƯƠNG 1 - Rubric: Phát biểu bài toán
\chapter{Bài toán dự báo sự kiện và ý tưởng Future-Guided Learning}
\label{chuong:baitoan}

\section{Đánh đổi giữa tính hữu dụng và độ chính xác trong dự báo sự kiện}

Dự báo sự kiện theo chuỗi thời gian (time-series event forecasting) là bài toán xác định
liệu một biến cố quan tâm --- chẳng hạn một cơn động kinh, một lần hỏng hóc thiết bị, hay một
thay đổi mạnh trong động lực của một hệ --- có xảy ra trong một khoảng thời gian tới hay
không, dựa trên các quan sát đã có. Bài toán này, cùng với bài toán dự báo hồi quy giá trị
tương lai của chuỗi, có một đặc điểm chung gây khó khăn cho việc thiết kế mô hình: mối quan hệ
ngược giữa thời điểm dự báo và chất lượng dự báo.

Nếu dự báo được thực hiện \textit{sớm}, tức còn cách xa thời điểm sự kiện xảy ra, thì kết quả
dự báo \textit{hữu dụng} vì còn đủ thời gian để phản ứng (ví dụ cảnh báo bệnh nhân hoặc bác sĩ
trước khi cơn động kinh xảy ra), nhưng lại \textit{khó và kém chính xác}, vì các tín hiệu báo
trước sự kiện còn mơ hồ, lẫn trong nhiễu của hoạt động bình thường. Ngược lại, nếu dự báo được
thực hiện \textit{muộn}, gần sát hoặc ngay tại thời điểm sự kiện, thì kết quả \textit{dễ và
chính xác} hơn nhiều, vì các dấu hiệu của sự kiện đã rõ ràng, nhưng lại \textit{không còn giá
trị thực tiễn}, vì không còn thời gian để hành động. Với bài toán dự báo hồi quy một chuỗi
thời gian xa về phía trước ($H$ bước), hiện tượng tương tự xảy ra: dự báo xa (tầm dự báo $H$
lớn) khó và sai số cao hơn dự báo gần (tầm dự báo nhỏ), vì bất định tích lũy theo thời gian,
đặc biệt với các hệ có tính hỗn loạn (chaotic).

Phần lớn các phương pháp học sâu cho dự báo chuỗi thời gian huấn luyện một mô hình duy nhất
trực tiếp trên nhiệm vụ khó (dự báo sớm, hoặc dự báo xa), buộc mô hình phải tự học cách xử lý
độ khó này mà không có thêm thông tin hỗ trợ nào khác ngoài chính dữ liệu đầu vào và nhãn thật.
Câu hỏi đặt ra là: liệu có thể tận dụng được một mô hình \textit{khác}, đã học tốt trên một
phiên bản \textit{dễ hơn} của cùng bài toán, để hỗ trợ việc huấn luyện mô hình chính giải
quyết phiên bản \textit{khó}?

\section{Ý tưởng Future-Guided Learning}

Future-Guided Learning (FGL), phương pháp đề xuất trong bài báo của Gunasekaran và cộng sự
\cite{gunasekaran2025fgl}, trả lời câu hỏi trên bằng cách tách bài toán thành hai mô hình:

\begin{itemize}
	\item \textbf{Teacher}: được huấn luyện trên một phiên bản \textit{dễ hơn} của bài toán,
	sử dụng dữ liệu ở \textit{tương lai gần hơn} so với student --- gần hoặc ngay tại thời
	điểm sự kiện xảy ra (với bài toán phân loại sự kiện), hoặc chỉ cần ngoại suy một bước thay
	vì nhiều bước (với bài toán hồi quy). Vì nhiệm vụ dễ, teacher học nhanh và đạt độ chính xác
	cao.
	\item \textbf{Student}: giải quyết nhiệm vụ \textit{khó} thực sự cần dùng trong thực tế ---
	dự báo sớm, hoặc dự báo xa --- chỉ dùng dữ liệu ở \textit{quá khứ tương đối} so với teacher.
	\item \textbf{Cầu nối}: trong lúc huấn luyện student, cùng một batch dữ liệu đầu vào được
	đưa qua cả teacher và student, rồi phân phối xác suất đầu ra của student được ép
	\textit{bắt chước} phân phối của teacher thông qua chưng cất tri thức (knowledge
	distillation), đo bằng độ phân kỳ Kullback--Leibler (KL).
\end{itemize}

Tác giả cho biết ý tưởng này lấy cảm hứng từ lý thuyết \textit{predictive coding} trong khoa
học thần kinh, theo đó não bộ liên tục tạo dự báo trước cho tín hiệu cảm giác sẽ đến, rồi điều
chỉnh dựa trên chênh lệch giữa dự báo và tín hiệu thực tế nhận được \cite{raoballard1999predictive}.
Trong FGL, vai trò ``tín hiệu sửa sai'' này được đóng bởi hàm mất mát chưng cất tri thức giữa
teacher và student: teacher đóng vai trò nguồn thông tin về tương lai gần, còn quá trình chưng
cất đóng vai trò cơ chế truyền thông tin đó về cho student, dù student không được phép ``nhìn
thấy'' trực tiếp dữ liệu tương lai mà teacher có.

Điểm mấu chốt cần làm rõ khi phát biểu bài toán là: cơ chế ``nhìn thấy tương lai'' của teacher
\textit{không} có nghĩa là teacher được truy cập dữ liệu mà student không được phép dùng khi
suy luận (inference) trên dữ liệu mới. Teacher chỉ được huấn luyện \textit{tách biệt}, trên
một nhiệm vụ khác (dễ hơn) với nhiệm vụ của student; khi student dự báo một chuỗi mới trong
thực tế, chỉ có student (đã được chưng cất tri thức từ teacher trong lúc huấn luyện) mới được
dùng để suy luận, teacher không còn tham gia.

\section{Phát biểu hình thức}

Xét một mô hình teacher $f_T$ và một mô hình student $f_S$, cùng xử lý chuỗi quan sát đầu vào
$\mathbf{x}$, nhưng được huấn luyện trên hai phiên bản nhãn khác nhau của cùng một đích dự
báo: $y^{(T)}$ (nhãn của nhiệm vụ dễ, teacher) và $y^{(S)}$ (nhãn của nhiệm vụ khó, student),
với $y^{(T)}$ và $y^{(S)}$ cùng nhằm tới một mốc đích --- hoặc cùng một thời điểm tuyệt đối
trong bài toán hồi quy, hoặc cùng một sự kiện trong bài toán phân loại --- nhưng khác nhau về
lượng thông tin ``tương lai'' mà đầu vào tương ứng của mỗi mô hình chứa.

Teacher được huấn luyện độc lập trước, tối thiểu hóa một hàm mất mát chuẩn (ví dụ sai số bình
phương trung bình với hồi quy, hoặc entropy chéo với phân loại) trên nhiệm vụ của nó. Sau khi
teacher đã huấn luyện xong, student được huấn luyện với hàm mất mát tổng hợp
\begin{equation}
	\mathcal{L}_{\text{student}} = \alpha \cdot \mathcal{L}_{\text{CE}}\!\left(y^{(S)}, f_S(\mathbf{x})\right) \;+\; (1-\alpha) \cdot T^2 \cdot D_{\mathrm{KL}}\!\left(\operatorname{softmax}\!\left(\frac{f_T(\mathbf{x})}{T}\right) \,\middle\|\, \operatorname{softmax}\!\left(\frac{f_S(\mathbf{x})}{T}\right)\right),
	\label{eq:fgl-loss}
\end{equation}
trong đó $f_T(\mathbf{x})$ và $f_S(\mathbf{x})$ là logits (đầu ra trước softmax) của teacher
và student trên \textit{cùng} đầu vào $\mathbf{x}$ của student, $\alpha \in [0,1]$ cân bằng
giữa việc tin vào nhãn thật và tin vào teacher, còn $T$ (temperature) làm mềm phân phối xác
suất trước khi so sánh. Chi tiết và ý nghĩa của từng tham số được trình bày trong
Chương~\ref{chuong:phuongphap}.

Cần nhấn mạnh rằng công thức \eqref{eq:fgl-loss} chỉ xác định được \textit{dạng toán học} của
cơ chế chưng cất; việc teacher và student \textit{cụ thể} nhìn vào phần dữ liệu nào, và nhãn
$y^{(T)}$, $y^{(S)}$ được định nghĩa ra sao, mới là phần quyết định ``tương lai'' nghĩa là gì
trong từng bài toán --- đây là điểm khác biệt chính giữa hai ứng dụng được tái hiện trong báo
cáo này (Chương~\ref{chuong:phuongphap}, mục~\ref{sec:cochetuonglai}).

\section{Khó khăn và thách thức}

Việc hiện thực hóa ý tưởng FGL gặp một số khó khăn chính sau.

Thứ nhất là việc áp dụng chưng cất tri thức --- vốn được định nghĩa trên phân phối xác suất
rời rạc của bài toán phân loại \cite{hinton2015distillation} --- cho bài toán \textit{hồi quy}
(dự báo giá trị liên tục của chuỗi thời gian). Cần một cách chuyển bài toán hồi quy thành bài
toán phân loại để có thể tính KL divergence giữa teacher và student.

Thứ hai là việc xác định cụ thể ``tương lai gần hơn'' nghĩa là gì cho từng loại bài toán. Với
bài toán hồi quy trên một chuỗi liên tục, không có khái niệm ``nhãn khác'' rõ ràng như bài toán
phân loại sự kiện; cần một cách lệch cửa sổ dữ liệu đầu vào sao cho teacher và student cùng
hướng tới một đích, nhưng xuất phát điểm khác nhau.

Thứ ba, với bài toán dự báo sự kiện y tế như cơn động kinh, dữ liệu ở giai đoạn ngay trước sự
kiện (tiền triệu, preictal) thường rất hiếm so với dữ liệu bình thường (interictal), gây mất
cân bằng lớp nghiêm trọng, ảnh hưởng đến cả việc huấn luyện teacher (nếu teacher cũng cần nhãn
sự kiện) và student.

Thứ tư là chi phí tính toán và độ phức tạp triển khai: FGL yêu cầu huấn luyện và lưu trữ hai
mô hình thay vì một, cùng với việc phải chạy suy luận của teacher trong suốt quá trình huấn
luyện student để tính hàm mất mát chưng cất.

\section{Đóng góp của bài báo gốc}

Từ các phân tích trên, đóng góp chính của bài báo gốc có thể tóm lược như sau
\cite{gunasekaran2025fgl}:
\begin{itemize}
	\item Đề xuất khung Future-Guided Learning, áp dụng được cho cả bài toán dự báo hồi quy và
	bài toán dự báo sự kiện nhị phân, chỉ bằng cách thay đổi cách định nghĩa dữ liệu và nhãn
	của teacher.
	\item Đề xuất cách rời rạc hóa bài toán hồi quy (chia khoảng giá trị thành các lớp) để áp
	dụng được chưng cất tri thức dựa trên KL divergence, vốn chỉ định nghĩa tự nhiên trên bài
	toán phân loại.
	\item Thực nghiệm trên ba bộ dữ liệu thuộc hai lĩnh vực khác nhau: chuỗi hỗn loạn tổng hợp
	Mackey--Glass (dự báo hồi quy), và hai bộ dữ liệu điện não đồ CHB-MIT và AES (dự báo cơn
	động kinh), báo cáo mức giảm 23,4\% sai số MSE trên Mackey--Glass và mức tăng 44,8\% AUC-ROC
	trung bình trên CHB-MIT so với mô hình không dùng FGL.
\end{itemize}
